\begin{thebibliography}{99}\footnotesize
\setlength{\itemsep}{0pt}\setlength{\parskip}{0pt}

\bibitem{bader2024} M.~Bader, M.~Abdelwanis, M.~Maalouf et H.~F.~Jelinek, « Detecting depression severity using weighted random forest and oxidative stress biomarkers », \emph{Sci. Rep.}, vol.~14, art.~16328, 2024, doi:10.1038/s41598-024-67251-y.

\bibitem{breiman2001} L.~Breiman, « Random forests », \emph{Mach. Learn.}, vol.~45, no~1, pp.~5--32, 2001.

\bibitem{chen2004} C.~Chen, A.~Liaw et L.~Breiman, « Using random forest to learn imbalanced data », Dept. of Statistics, Univ. of California, Berkeley, Tech. Rep.~666, 2004.

\bibitem{maalouf2014} M.~Maalouf et M.~Siddiqi, « Weighted logistic regression for large-scale imbalanced and rare events data », \emph{Knowl.-Based Syst.}, vol.~59, pp.~142--148, 2014.

\bibitem{zhang2016} C.~Zhang, Y.~Li, Z.~Yu et F.~Tian, « A weighted random forest approach to improve predictive performance for power system transient stability assessment », in \emph{Proc. IEEE PES APPEEC}, 2016, pp.~1259--1263.

\bibitem{chawla2002} N.~V.~Chawla, K.~W.~Bowyer, L.~O.~Hall et W.~P.~Kegelmeyer, « SMOTE: Synthetic minority over-sampling technique », \emph{J. Artif. Intell. Res.}, vol.~16, pp.~321--357, 2002.

\bibitem{bergstra2012} J.~Bergstra et Y.~Bengio, « Random search for hyper-parameter optimization », \emph{J. Mach. Learn. Res.}, vol.~13, pp.~281--305, 2012.

\bibitem{bergstra2011} J.~Bergstra, R.~Bardenet, Y.~Bengio et B.~K\'egl, « Algorithms for hyper-parameter optimization », in \emph{Adv. Neural Inf. Process. Syst. (NIPS)}, vol.~24, 2011, pp.~2546--2554.

\bibitem{akiba2019} T.~Akiba, S.~Sano, T.~Yanase, T.~Ohta et M.~Koyama, « Optuna: A next-generation hyperparameter optimization framework », in \emph{Proc. ACM SIGKDD}, 2019, pp.~2623--2631.

\bibitem{kohavi1995} R.~Kohavi, « A study of cross-validation and bootstrap for accuracy estimation and model selection », in \emph{Proc. IJCAI}, 1995, pp.~1137--1143.

\bibitem{cawley2010} G.~C.~Cawley et N.~L.~C.~Talbot, « On over-fitting in model selection and subsequent selection bias in performance evaluation », \emph{J. Mach. Learn. Res.}, vol.~11, pp.~2079--2107, 2010.

\bibitem{kaufman2012} S.~Kaufman, S.~Rosset, C.~Perlich et O.~Stitelman, « Leakage in data mining: Formulation, detection, and avoidance », \emph{ACM Trans. Knowl. Discov. Data}, vol.~6, no~4, art.~15, 2012.

\bibitem{nadeau2003} C.~Nadeau et Y.~Bengio, « Inference for the generalization error », \emph{Mach. Learn.}, vol.~52, no~3, pp.~239--281, 2003.

\bibitem{dietterich1998} T.~G.~Dietterich, « Approximate statistical tests for comparing supervised classification learning algorithms », \emph{Neural Comput.}, vol.~10, no~7, pp.~1895--1923, 1998.

\bibitem{demsar2006} J.~Dem\v{s}ar, « Statistical comparisons of classifiers over multiple data sets », \emph{J. Mach. Learn. Res.}, vol.~7, pp.~1--30, 2006.

\bibitem{friedman1937} M.~Friedman, « The use of ranks to avoid the assumption of normality implicit in the analysis of variance », \emph{J. Am. Stat. Assoc.}, vol.~32, no~200, pp.~675--701, 1937.

\bibitem{conover1999} W.~J.~Conover, \emph{Practical Nonparametric Statistics}, 3rd ed. New York: Wiley, 1999.

\bibitem{pereira2015} D.~G.~Pereira, A.~Afonso et F.~M.~Medeiros, « Overview of Friedman's test and post-hoc analysis », \emph{Commun. Stat. Simul. Comput.}, vol.~44, no~10, pp.~2636--2653, 2015.

\bibitem{wilcoxon1945} F.~Wilcoxon, « Individual comparisons by ranking methods », \emph{Biometrics Bull.}, vol.~1, no~6, pp.~80--83, 1945.

\bibitem{kroenke2001} K.~Kroenke, R.~L.~Spitzer et J.~B.~W.~Williams, « The PHQ-9: Validity of a brief depression severity measure », \emph{J. Gen. Intern. Med.}, vol.~16, no~9, pp.~606--613, 2001.

\bibitem{altmann2010} A.~Altmann, L.~Tolo\c{s}i, O.~Sander et T.~Lengauer, « Permutation importance: a corrected feature importance measure », \emph{Bioinformatics}, vol.~26, no~10, pp.~1340--1347, 2010.

\bibitem{saito2015} T.~Saito et M.~Rehmsmeier, « The precision-recall plot is more informative than the ROC plot when evaluating binary classifiers on imbalanced datasets », \emph{PLoS ONE}, vol.~10, no~3, e0118432, 2015.

\bibitem{fawcett2006} T.~Fawcett, « An introduction to ROC analysis », \emph{Pattern Recognit. Lett.}, vol.~27, no~8, pp.~861--874, 2006.

\bibitem{hastie2009} T.~Hastie, R.~Tibshirani et J.~Friedman, \emph{The Elements of Statistical Learning}, 2nd ed. New York: Springer, 2009.

\bibitem{sau2017} A.~Sau et I.~Bhakta, « Predicting anxiety and depression in elderly patients using machine learning technology », \emph{Healthc. Technol. Lett.}, vol.~4, no~6, pp.~238--243, 2017.

\bibitem{priya2020} A.~Priya, S.~Garg et N.~P.~Tigga, « Predicting anxiety, depression and stress in modern life using machine learning algorithms », \emph{Procedia Comput. Sci.}, vol.~167, pp.~1258--1267, 2020.

\bibitem{zulfiker2021} M.~S.~Zulfiker, N.~Kabir, A.~A.~Biswas, T.~Nazneen et M.~S.~Uddin, « An in-depth analysis of machine learning approaches to predict depression », \emph{Curr. Res. Behav. Sci.}, vol.~2, art.~100044, 2021.

\bibitem{hidaka2012} B.~H.~Hidaka, « Depression as a disease of modernity: Explanations for increasing prevalence », \emph{J. Affect. Disord.}, vol.~140, pp.~205--214, 2012.

\bibitem{black2015} C.~N.~Black, M.~Bot, P.~G.~Scheffer, P.~Cuijpers et B.~W.~J.~H.~Penninx, « Is depression associated with increased oxidative stress? A systematic review and meta-analysis », \emph{Psychoneuroendocrinology}, vol.~51, pp.~164--175, 2015.

\bibitem{jelinek2014} H.~F.~Jelinek, K.~Yatsko, A.~Stranieri et S.~Venkatraman, « Novel data mining techniques for incomplete clinical data in diabetes management », \emph{Br. J. Appl. Sci. Technol.}, vol.~4, pp.~4591--4606, 2014.

\bibitem{lee2015} C.~Lee \emph{et al.}, « The mitochondrial-derived peptide MOTS-c promotes metabolic homeostasis and reduces obesity and insulin resistance », \emph{Cell Metab.}, vol.~21, no~3, pp.~443--454, 2015.

\bibitem{migliaccio1999} E.~Migliaccio \emph{et al.}, « The p66\textsuperscript{shc} adaptor protein controls oxidative stress response and life span in mammals », \emph{Nature}, vol.~402, pp.~309--313, 1999.

\bibitem{pedregosa2011} F.~Pedregosa \emph{et al.}, « Scikit-learn: Machine learning in Python », \emph{J. Mach. Learn. Res.}, vol.~12, pp.~2825--2830, 2011.

\bibitem{lemaitre2017} G.~Lema\^itre, F.~Nogueira et C.~K.~Aridas, « Imbalanced-learn: A Python toolbox to tackle the curse of imbalanced datasets in machine learning », \emph{J. Mach. Learn. Res.}, vol.~18, no~17, pp.~1--5, 2017.

\bibitem{terpilowski2019} M.~Terpilowski, « scikit-posthocs: Pairwise multiple comparison tests in Python », \emph{J. Open Source Softw.}, vol.~4, no~36, art.~1169, 2019.

\bibitem{harris2020} C.~R.~Harris \emph{et al.}, « Array programming with NumPy », \emph{Nature}, vol.~585, pp.~357--362, 2020.

\bibitem{virtanen2020} P.~Virtanen \emph{et al.}, « SciPy 1.0: Fundamental algorithms for scientific computing in Python », \emph{Nat. Methods}, vol.~17, pp.~261--272, 2020.

\bibitem{mckinney2010} W.~McKinney, « Data structures for statistical computing in Python », in \emph{Proc. 9th Python in Science Conf.}, 2010, pp.~56--61.

\bibitem{hunter2007} J.~D.~Hunter, « Matplotlib: A 2D graphics environment », \emph{Comput. Sci. Eng.}, vol.~9, no~3, pp.~90--95, 2007.

\end{thebibliography}
